% META: {"id": "TCOMPL-ATT-CO-037", "chapitre": "TCOMPL-TEMPS-ATTENTE", "type_objet": "corrige", "exercice_ref": "TCOMPL-ATT-EX-037", "capacites_codes": ["C7"], "sources_inspiration": [], "status": "needs_review"}
% BEGIN-VERIFY
% from sympy import symbols, integrate, Rational, simplify
% x = symbols('x', real=True)
% f = Rational(1, 15)
% assert integrate(f, (x, 0, 15)) == 1
% assert integrate(x * f, (x, 0, 15)) == Rational(15, 2)
% variance = integrate(x**2 * f, (x, 0, 15)) - Rational(15, 2)**2
% assert variance == Rational(75, 4)
% assert integrate(f, (x, 0, 5)) == integrate(f, (x, 10, 15))
% END-VERIFY
\begin{corrige}{TCOMPL-ATT-EX-037}
\textbf{1.} La densité est constante sur $[0\,;15]$ : $f(x) = \dfrac{1}{15-0} = \dfrac{1}{15}$ pour $x \in [0\,;15]$, et $f(x)=0$ ailleurs. Elle est positive, et
\[
  \int_0^{15} \frac{1}{15}\,\mathrm{d}x = \frac{1}{15}\times 15 = 1 .
\]
C'est bien une densité de probabilité.

\textbf{2.} Pour $x \in [0\,;15]$ :
\[
  F(x) = P(X\leq x) = \int_0^{x} \frac{1}{15}\,\mathrm{d}t = \frac{x}{15}.
\]
On a bien $F(0)=0$ et $F(15)=1$.

\textbf{3.} $P(X\leq5) = \dfrac{5}{15} = \dfrac{1}{3}$ et $P(10\leq X\leq 15) = \dfrac{15-10}{15} = \dfrac{1}{3}$. Les deux probabilités sont égales : pour une loi uniforme, seule la \emph{longueur} de l'intervalle compte, pas sa position.

\textbf{4.} $E(X) = \dfrac{0+15}{2} = 7{,}5$ minutes : c'est l'attente moyenne d'un voyageur arrivant au hasard, soit la moitié de l'intervalle entre deux métros. Et
\[
  V(X) = \frac{(15-0)^2}{12} = \frac{225}{12} = 18{,}75 \ \text{(minutes}^2\text{)},
\]
soit un écart-type d'environ $4{,}33$ minutes.
\end{corrige}
